\section{Further research questions and concrete next targets}
\label{clc:sec:questions}

The following questions concern further exact identities rather than
inequalities or asymptotic error bounds. They also distinguish genuine
remaining problems from the two continuation questions resolved here.

\begin{enumerate}[label=\textbf{Q\arabic*.},leftmargin=2.8em,itemsep=9pt]
\item \textbf{A compact multivariate tilt generating function.}
For the constraint $\sum_jx_j=\mu$, introduce independent factors
$e^{\lambda_jx_j}$. Their transform becomes
$\prod_jg_{q_j}(p-\lambda_j)(p+A-\lambda_j)^{-s_j}$.
Find a compact representation in a specified multivariate Lerch or
Lauricella family, with every collision and normalization explicit,
that avoids an outer center-expansion series. The present moment theorem
already determines every Taylor coefficient, but does not supply such
a compact resummation. A characterization of its genuinely finite
sectors would be more informative than an unqualified ``closed form.''

\item \textbf{The complementary negative-order parity sector.}
Theorem~\ref{clc:thm:parity} evaluates the sector $r+N$ odd at $A=0$.
For $r+N$ even and incommensurate scales, seek exact representations in
Barnes-type or multiple Hurwitz functions that retain the full scale
symmetry. A correct formula should reduce to the finite commensurate
closure when all scale ratios become rational, without assuming that
the limiting pole multiplicities are unchanged.

\item \textbf{Minimal cyclotomic coordinates and conductor descent.}
The finite tables live in explicit real cyclotomic algebras, but those
algebras and the Hurwitz argument grids are often redundant. Determine
an exact reduction algorithm for the coefficient field and the
Hurwitz-coordinate conductor simultaneously. The target is a canonical
functional reduction with exact certificates, not numerical independence
of a chosen set of periods. The two-scale beta and character-three
formulas are required base cases.

\item \textbf{Finite closure without pure order addition.}
Alignment is characterized exactly for pure addition. It does not
classify all finite identities with unequal centers. Positive integer
orders permit finite partial fractions, while our arbitrary-order
series is usually infinite. Classify exceptional affine-center and
order subvarieties on which the outer series terminates or resums
finitely. Nonpositive integer orders, repeated centers, and rational
center differences provide distinct test cases and should not be
conflated.

\item \textbf{Beyond a single common-step elementary kernel class.}
Theorem~\ref{clc:thm:commensurate} allows one shared positive $Q$.
Determine the exact obstruction when finitely many different denominator
periods are allowed in the elementary representation. The uncancelled
pole support alone may no longer suffice; the Laurent coefficient
sequences along intersecting progressions carry additional information.
Any stronger theorem should state its representational class precisely,
and should remain compatible with the universal isolated parity values.

\item \textbf{Several independent multiplicative constraints.}
Replace the single linear condition on logarithms by an integer matrix
of constraints. The transform is then a product of cosecants and affine
resolvents on a higher-dimensional space. Seek a finite exact
polylogarithm/Hurwitz closure when the pole hyperplane arrangement is
rational and the resolvents are aligned along the appropriate dual
subspaces. A first nontrivial target is two independent constraints
with four or five factors, including intersecting-pole residues and all
independent order derivatives.

\item \textbf{Colored profiles and exact contour-crossing terms.}
Replace the negative argument of the Lerch profile by a root-of-unity
multiple, and allow complex logarithmic shifts. Derive the correct
weighted strip and the finite residue or branch terms when a pole
crosses the chosen contour. The jump cancellation of
Theorem~\ref{clc:thm:finite} is a special aligned case, not a general
rule that all colored jumps vanish. Exact quarter-turn specializations
could connect this extension to the manuscript's Gaussian identities.

\item \textbf{A canonical finite basis for mixed moments.}
The algorithm in Theorem~\ref{clc:thm:moments} is terminating but not
claimed minimal. Determine all functional syzygies among its differentiated
Laurent tables and produce a canonical reduced representative that
respects center differentiation, scale covariance, and
\eqref{clc:eq:moment-sum}. The equal-scale central-factorial recurrence
is an instructive subcase; unequal periods introduce new redundancies.

\item \textbf{Machine-checked analytic infrastructure.}
Formalize weighted entire families, bilateral convolution inversion,
and the finite period-residue argument in a proof assistant. The exact
coefficient tables and negative controls can be used as finite fixtures,
but they do not replace a formal proof of contour deformation, normal
convergence, or branch cancellation. The first useful formal milestone
is the all-complex-order two-scale theorem with its removable resonances.
\end{enumerate}

The commensurate-dilation problem and the unequal-shift order-jet
expansion posed in \cite{RLC} now have complete proofs in the stated
settings. The next advances should go beyond those solved settings,
while preserving their normalization and convergence checks. In
particular, connecting these integral families to the unresolved
canonical Gaussian reduction vectors would require a separate explicit
identity or exact reduction certificate; it is not a consequence of
finite closure alone.
